% Generated draft: requires mathematical and source review.
\documentclass[11pt]{article}
\usepackage[margin=25mm]{geometry}
\usepackage{fontspec}
\setmainfont{Latin Modern Roman}
\usepackage{amsmath,amssymb,mathtools,unicode-math}
\setmathfont{Latin Modern Math}
\usepackage{enumitem,needspace,xurl,hyperref}
\setlength{\parindent}{0pt}
\setlength{\parskip}{5pt}
\setlength{\emergencystretch}{3em}
\newcommand{\N}{\mathbb{N}}
\newcommand{\Z}{\mathbb{Z}}
\newcommand{\Q}{\mathbb{Q}}
\newcommand{\R}{\mathbb{R}}
\newcommand{\C}{\mathbb{C}}
\newcommand{\F}{\mathbb{F}}
\newcommand{\A}{\mathbb{A}}
\newcommand{\PP}{\mathbb{P}}
\newcommand{\E}{\mathbb{E}}
\newcommand{\eps}{\varepsilon}
\newcommand{\dd}{\,\mathrm{d}}
\begin{document}
\section*{UnsolvedMath ID: 30006596}
\textbf{Source problem code:} catalogue entry
\subsection{Definitions and mathematical statement}
\textbf{Problem (attributed in the supplied source transcription to Kam Cheong Au).} Construct a \(q\)-analogue of Ramanujan's \(9801\) series for \(1/\pi\) that exhibits modular properties. The cleaned registry statement also permits a mock-modular analogue.

Here \(q\) denotes a deformation parameter, but its domain is not specified in the supplied record. The expressions ``\(q\)-analogue,'' ``modular,'' and ``mock-modular'' are requirements whose precise meanings remain to be specified for this problem: the record gives no required specialization or limiting procedure recovering the classical series, no admissible class of identities, and no modular group, weight, multiplier, transformation law, or mock-modular completion data.

The original question asks more broadly for such analogues of some of Ramanujan's \(1/\pi\) series and presents the \(9801\) series as an example; the cleaned statement selects that example as the target. The supplied transcription of the displayed classical formula is corrupted, so its exact summand and normalization are not reproduced here by conjectural reconstruction. It also contains an incomplete comparison with a simpler \(1/\pi\) identity, whose formula and subsequent discussion cannot be recovered reliably from the supplied text. Consequently, this is a statement of the requested construction with explicit unresolved qualifications, rather than a fully specified existence proposition.
\subsection{Short English statement}
Find a \(q\)-deformation of Ramanujan's \(9801\) formula for \(1/\pi\) with modular properties; the cleaned record also allows mock-modular properties. The exact classical formula, recovery condition, and required transformation properties need clarification from the cited report before the target can be made fully precise.
\subsection{Sources}
\begin{itemize}
\item Local UnsolvedMath registry, \texttt{data/unsolvedmath/problems.json}, permanent ID \(30006596\), source problem code \texttt{catalogue entry}, title: Modular \(q\)-Analogues of Ramanujan's \(9801\) Series. The supplied original transcription attributes the proposal to Kam Cheong Au, although the structured proposer field is empty.
\item \emph{Multiple Zeta Values and Modular Forms} (2026), Oberwolfach Reports, Report 7/2026, pp.\ 433--486, DOI \texttt{10.4171/owr/2026/7}: \url{https://doi.org/10.4171/owr/2026/7}. This is the supplied report citation; the record supplies no page locating the individual problem.
\item Campbell, \(q\)-analogues of pi-formulas, supplied literature-source label: \url{https://arxiv.org/abs/2509.05897}.
\item Guo--Zudilin, Ramanujan-type \(q\)-analogues, supplied literature-source label: \url{https://doi.org/10.1080/10652469.2018.1454448}.
\item Supplied literature assessment dated \texttt{2026-08-22}: modular or mock-modular \(q\)-analogues are reported for several other Ramanujan formulas, while the cited literature is reported not to establish one for the displayed \(9801\) series; the problem is labeled open. This assessment is supplied dated curation derived from the Oberwolfach record, not independently verified here. The supplied background expressly requests independent verification of time-sensitive status and cited claims.
\end{itemize}
\end{document}
